\section{Experiments and Results}
\label{sec:experiments}

Three groups of experiments evaluate the proposed swing-arm CPP method. The first compares it with Fields2Cover to clarify the difference between active arm sweeping and fixed-width strip-based CPP. The second compares rectangular-candidate scan strategies to determine how decomposition affects coverage quality and redundancy in complex spaces. The third varies base speed at a fixed arm angular velocity to expose how velocity matching affects coverage and execution time.

All three groups use two base maps. \texttt{map\_test1} contains four regular rectangular obstacles and tests basic obstacle avoidance and coverage. \texttt{parkingGraph} contains densely distributed small obstacles and better resembles a construction floor partitioned by columns, openings, and local structures. The scan-strategy experiment additionally uses three directional maps:
\begin{itemize}[nosep]
  \item \texttt{scan\_comb\_corridors}: horizontal comb-like corridors;
  \item \texttt{scan\_mixed\_weave}: interwoven horizontal and vertical structures; and
  \item \texttt{scan\_offset\_rooms}: rooms connected with offsets.
\end{itemize}
These structures distinguish the effects of row-only, column-only, and bidirectional scanning.

Every map is a binary occupancy image in which white pixels are free target area and black pixels are obstacles or unavailable regions. The resolution is $0.5$~m/pixel. Except in the speed experiment, the same parameters are used throughout: disc radius $r_d=0.75$~m, arm length $L=1.20$~m, base width $0.80$~m, maximum base linear velocity $2.00$~m/s, maximum base angular velocity $30^\circ$/s, and maximum arm angular velocity $90^\circ$/s. The minimum rectangle side length is $2.5$~m, the greedy-selection weights are $\alpha=0.5$ and $\beta=2.0$, and the random seed is 42.

\subsection{Comparison with Fields2Cover}

Fields2Cover is an appropriate fixed-width baseline and a representative open-source framework for this class of problem. It is supported by a peer-reviewed publication and continuing open-source development, integrates multiple coverage-planning algorithms and objectives, and includes headland generation, swath generation, route ordering, and curvature-constrained path connections. It supports non-convex regions, internal obstacles, fixed working and vehicle widths, turning radii, and speed parameters. Its functional completeness, algorithmic integration, and reproducibility make it a strong contemporary baseline for fixed-width CPP \cite{mier2023,fields2cover}. The comparison is not intended to rank implementation details, but to expose differences in coverage geometry and executability between an active swing-arm model and an integrated, fixed-width strip-planning framework in obstacle-rich spaces.

For Fields2Cover, the working width is $1.50$~m, equal to the diameter of the proposed method's $0.75$~m-radius disc. With base width $0.80$~m and safety margin $0.05$~m, the original free region is eroded by $0.45$~m. Both methods are evaluated over their complete trajectories. For the proposed method, coverage is the area swept by the disc along the base--arm trajectory; for Fields2Cover, it is the fixed-width tool sweep along work, turning, and transfer segments. To align evaluation with the actual planning domain, the Fields2Cover denominator is the eroded free-space area. Its fixed-width sweep uses flat strip ends and mitered path corners. The work-strip speed is $2.0$~m/s, while turns and connecting paths use $0.5$~m/s.

\begin{table}[htbp]
\centering
\caption{Comparison between the proposed method and Fields2Cover.}
\label{tab:baseline-comparison}
\scriptsize
\setlength{\tabcolsep}{3pt}
\renewcommand{\arraystretch}{1.18}
\begin{tabular}{@{}p{1.95cm}p{1.45cm}rrrrp{2.55cm}@{}}
\toprule
Map & Method & \shortstack{Evaluation\\area (m$^2$)} & \shortstack{Covered\\area (m$^2$)} & \shortstack{Coverage\\ratio (\%)} & \shortstack{Execution/planning\\time (s/s)} & Path check \\
\midrule
\texttt{map\_test1} & Proposed & 611.00 & 575.00 & 94.11 & 854.73/1.298 & Passed \\
 & Fields2Cover & 533.17 & 418.00 & 78.40 & 295.50/0.075 & 4.51~m$^2$ outside \\
\texttt{parking}\allowbreak\texttt{Graph} & Proposed & 4294.00 & 4110.50 & 95.73 & 5844.73/17.981 & Passed \\
 & Fields2Cover & 4046.79 & 3756.85 & 92.84 & 2883.81/6.339 & 20.00~m$^2$ outside \\
\bottomrule
\end{tabular}
\end{table}

The evaluated areas in Table~\ref{tab:baseline-comparison} differ between methods, so their coverage ratios should not be interpreted as same-domain percentage-point differences. The Fields2Cover routes fail the path check; their execution and planning times are reported only for the generated outputs.

On \texttt{map\_test1}, Fields2Cover has shorter generation and nominal execution times with a $1.50$~m fixed width equal to the disc diameter. Its complete trajectory sweeps 418.00~m$^2$ (78.40\%), compared with 575.00~m$^2$ (94.11\%) for the proposed method. Its connecting segments, however, sweep 4.51~m$^2$ outside the admissible region. The recorded 295.50~s execution time and 0.075~s planning time therefore describe an invalid generated path rather than an executable solution. This illustrates the principal benefit of the swing-arm configuration: the base need not drive along every strip that the disc must cover. Arm motion extends the covered region and creates a wider time-varying sweep within the same free space.

On \texttt{parkingGraph}, the proposed method covers 4110.50~m$^2$, or 95.73\% of its evaluation domain. Fields2Cover sweeps 3756.85~m$^2$ (92.84\%) along the complete trajectory. Yet its connecting path for the $1.50$~m tool leaves the free space, creating 20.00~m$^2$ of invalid base sweep; the corresponding 2883.81~s execution time and 6.339~s planning time are likewise not valid baseline scores. The problem is concentrated near dense local obstacles. Although individual swaths fit inside the eroded region, curvature-constrained connections between swaths require additional room. Without a broad transfer area analogous to an agricultural headland, these connections can cross the boundary of the original free space. Active arm sweeping partially decouples base travel from end-disc coverage and preserves greater coverage capability and path continuity in such complex regions.

\begin{figure}[htbp]
  \centering
  \includegraphics[width=0.96\linewidth]{figure_2_coverage_comparison_en.png}
  \caption{Coverage paths and swept regions for the proposed method and Fields2Cover. Panels (a) and (b) show \texttt{map\_test1} (labeled \texttt{map\_simple} inside the image), and panels (c) and (d) show \texttt{parkingGraph}; the proposed method is on the left and Fields2Cover on the right. Blue is the full-trajectory swept region, white or light areas are uncovered free space in the evaluation domain, black is obstacle or unavailable space, and red marks base collisions or invalid sweep. The Fields2Cover route fails the base-collision check on both maps: 4.51~m$^2$ lies outside the valid region on \texttt{map\_test1} and 20.00~m$^2$ on \texttt{parkingGraph}. Its displayed coverage and timing values therefore describe generated, non-executable routes.}
  \label{fig:baseline-comparison}
\end{figure}

\subsection{Effect of the Rectangle-Candidate Scan Strategy}

The preceding comparison shows that an active arm can create coverage patterns unlike those of fixed-width strips. Realizing this advantage consistently also requires a decomposition that represents free space appropriately. A candidate generator biased toward one scan direction can produce fragmented or heavily overlapping rectangles; even when the resulting trajectory is executable, unnecessary transfers may follow. This experiment traces how scan direction propagates through the rectangle set, visitation sequence, and complete trajectory.

The compared strategies are row-only, column-only, and row-then-column scanning. Robot parameters, minimum rectangle size, greedy weights, random seed, connectivity completion, visitation ordering, and local coverage generation are fixed. Measures include greedy coverage, final decomposition coverage, decomposition overlap, full-trajectory coverage, and total execution time, distinguishing representation quality from execution performance. Candidates are first generated and greedily selected from a binary map. After hole filling and connectivity completion, the same trajectory generation and disc-sweep evaluation method is applied in every case.

\begin{longtable}{@{}p{2.50cm}p{1.60cm}>{\raggedleft\arraybackslash}p{1.55cm}>{\raggedleft\arraybackslash}p{1.65cm}>{\raggedleft\arraybackslash}p{1.45cm}>{\raggedleft\arraybackslash}p{1.60cm}>{\raggedleft\arraybackslash}p{1.55cm}@{}}
\caption{Results of three candidate-scan strategies on five maps.}\label{tab:scan-strategy}\\
\toprule
Map & Strategy & \shortstack{Greedy\\coverage} & \shortstack{Final\\partition\\coverage} & \shortstack{Partition\\overlap} & \shortstack{Trajectory\\coverage} & \shortstack{Total time\\(s)} \\
\midrule
\endfirsthead
\toprule
Map & Strategy & \shortstack{Greedy\\coverage} & \shortstack{Final\\partition\\coverage} & \shortstack{Partition\\overlap} & \shortstack{Trajectory\\coverage} & \shortstack{Total time\\(s)} \\
\midrule
\endhead
\texttt{map\_test1} & row-only & 72.01\% & 100.00\% & 34.21\% & 93.04\% & 696.89 \\
 & col-only & 72.01\% & 100.00\% & 34.21\% & 93.04\% & 696.89 \\
 & row-then-col & 100.00\% & 100.00\% & 0.00\% & 94.11\% & 854.73 \\
\midrule
\texttt{parking}\allowbreak\texttt{Graph} & row-only & 87.32\% & 100.00\% & 75.34\% & 95.73\% & 5844.73 \\
 & col-only & 87.32\% & 100.00\% & 75.34\% & 95.90\% & 5896.03 \\
 & row-then-col & 87.32\% & 100.00\% & 75.34\% & 95.73\% & 5844.73 \\
\midrule
Comb corridors & row-only & 87.18\% & 100.00\% & 34.46\% & 93.82\% & 2065.09 \\
 & col-only & 86.53\% & 100.00\% & 31.46\% & 93.94\% & 2032.55 \\
 & row-then-col & 96.22\% & 100.00\% & 27.13\% & 94.01\% & 1986.42 \\
\midrule
Mixed weave & row-only & 84.23\% & 100.00\% & 47.89\% & 95.26\% & 1532.65 \\
 & col-only & 83.17\% & 100.00\% & 52.59\% & 95.49\% & 1575.63 \\
 & row-then-col & 96.76\% & 100.00\% & 40.74\% & 95.70\% & 1522.77 \\
\midrule
Offset rooms & row-only & 82.95\% & 100.00\% & 38.78\% & 91.97\% & 1496.73 \\
 & col-only & 81.81\% & 100.00\% & 36.74\% & 91.51\% & 1443.31 \\
 & row-then-col & 99.14\% & 100.00\% & 24.46\% & 92.06\% & 1374.38 \\
\midrule
Five-map mean & row-only & -- & -- & -- & 93.96\% & 2327.22 \\
 & col-only & -- & -- & -- & 93.98\% & 2328.88 \\
 & row-then-col & -- & -- & -- & 94.32\% & 2316.60 \\
\bottomrule
\end{longtable}

The two general maps first establish basic behavior in non-directional scenes. On \texttt{map\_test1}, row-then-column achieves the highest full-trajectory coverage (94.11\%) and reduces partition overlap to zero, but its 854.73~s execution time exceeds the 696.89~s of either one-direction strategy. In this regular obstacle scene, bidirectional scanning trades longer execution for higher coverage. On \texttt{parkingGraph}, row-then-column and row-only select the same rectangles and give identical results. Column-only provides slightly higher coverage (95.90\%), whereas row-only and row-then-column tie for the shortest total time (5844.73~s). Bidirectional scanning therefore does not dominate every metric in every scene. When free space has a pronounced principal direction or row candidates already capture the structure, a one-direction strategy may be faster or marginally improve coverage.

The three directional maps then isolate scan-direction bias; they are not intended to prove that bidirectional scanning is optimal on arbitrary maps. Horizontal corridors, orthogonal interweaving, and offset connections amplify residual regions created by a single scan direction. On all three maps, row-then-column gives the highest greedy coverage and lowest partition overlap. Column candidates supplement regions that row candidates represent poorly rather than merely increasing the rectangle count. This improvement propagates to execution: row-then-column gives both the highest full-trajectory coverage and shortest time on all three maps, respectively 94.01\%/1986.42~s, 95.70\%/1522.77~s, and 92.06\%/1374.38~s. Their mean is 93.93\%/1627.86~s. When a map truly contains residual structures in both directions, complementary row and column candidates reduce directional bias and redundant visitation while increasing useful sweep along the complete trajectory.

Across all five maps, row-then-column averages 94.32\% coverage, compared with 93.96\% for row-only and 93.98\% for column-only. Its mean execution time of 2316.60~s is also slightly below 2327.22~s and 2328.88~s. It is therefore best characterized as a robust default: it is more likely to preserve high coverage without a substantial time penalty when horizontal and vertical residual structures coexist. For a regular region with a single principal direction, a one-direction strategy may still be preferable under the specific coverage--schedule trade-off. Figure~\ref{fig:scan-strategy} visualizes the directional complementarity on \texttt{scan\_mixed\_weave}.

\begin{figure}[htbp]
  \centering
  \includegraphics[width=0.96\linewidth]{figure_3_segmentation_order_en.png}
  \caption{Candidate-scan strategies on \texttt{scan\_mixed\_weave}: (a) original map; (b) row-only decomposition; (c) column-only decomposition; and (d) row-then-column decomposition. The map contains narrow horizontal and vertical free regions. Single-direction candidate sets leave residual structures in the other orientation, whereas sequential row and column candidates form a more complete decomposition with less overlap.}
  \label{fig:scan-strategy}
\end{figure}

\subsection{Effect of Base--Arm Velocity Coupling on Coverage}

The first two experiments address route and decomposition quality. Swing-arm coverage also has an execution issue that is less prominent in fixed-width CPP: the same geometric route produces different disc-sweep densities at different base speeds. As the base advances faster, it travels farther between adjacent arm sweeps; slower travel improves coverage but substantially increases task time. Reporting a geometric path without analyzing its velocity relationship to the arm therefore cannot establish the coverage quality and time cost of actual execution.

This experiment tests the analysis in Section~3.5.3. At fixed arm angular velocity, increasing base speed should shorten execution but enlarge the spacing between adjacent sweeps and reduce coverage. Keeping the rectangle set, visitation order, and local-work structure fixed isolates the effect of base speed.

The arm angular velocity is fixed at $\omega_a=90^\circ$/s. The same planned solutions are executed on \texttt{map\_test1} and \texttt{parkingGraph} with base speed increasing from 0.5 to 3.0~m/s. A velocity ratio places linear and angular speed on a common scale:
\begin{equation*}
\rho=\frac{v}{\omega_a},
\end{equation*}
with units of m/degree. It is the base distance traveled while the arm rotates through $1^\circ$. A larger $\rho$ increases forward spacing between adjacent sweeps and makes coverage gaps more likely. Coverage, total time, and time per unit covered area are recorded to expose the quality--efficiency trade-off.

\begin{table}[htbp]
\centering
\caption{Full-trajectory coverage and execution time at different base speeds.}
\label{tab:speed-study}
\scriptsize
\setlength{\tabcolsep}{3pt}
\renewcommand{\arraystretch}{1.12}
\begin{tabular}{@{}p{2.15cm}>{\raggedleft\arraybackslash}p{1.35cm}>{\raggedleft\arraybackslash}p{1.7cm}>{\raggedleft\arraybackslash}p{1.9cm}>{\raggedleft\arraybackslash}p{1.7cm}>{\raggedleft\arraybackslash}p{2.3cm}@{}}
\toprule
Map & \shortstack{Base speed\\(m/s)} & \shortstack{$\rho$\\(m/degree)} & \shortstack{Full-path\\coverage} & \shortstack{Total time\\(s)} & \shortstack{Time per covered\\area (s/m$^2$)} \\
\midrule
\texttt{map\_test1} & 0.5 & 0.00556 & 96.44\% & 1680.45 & 2.852 \\
 & 1.0 & 0.01111 & 95.46\% & 1095.53 & 1.878 \\
 & 1.5 & 0.01667 & 88.83\% & 928.81 & 1.711 \\
 & 2.0 & 0.02222 & 88.75\% & 829.74 & 1.530 \\
 & 2.5 & 0.02778 & 86.25\% & 772.26 & 1.465 \\
 & 3.0 & 0.03333 & 85.27\% & 735.12 & 1.411 \\
\midrule
\texttt{parking}\allowbreak\texttt{Graph} & 0.5 & 0.00556 & 99.15\% & 14525.14 & 3.412 \\
 & 1.0 & 0.01111 & 96.65\% & 8346.21 & 2.011 \\
 & 1.5 & 0.01667 & 93.74\% & 6214.22 & 1.544 \\
 & 2.0 & 0.02222 & 92.33\% & 5130.64 & 1.294 \\
 & 2.5 & 0.02778 & 91.03\% & 4526.86 & 1.158 \\
 & 3.0 & 0.03333 & 90.32\% & 4120.11 & 1.062 \\
\bottomrule
\end{tabular}
\end{table}

Table~\ref{tab:speed-study} shows a clear coverage--time trade-off, although the full-trajectory results make it more moderate than an analysis of work strips alone might suggest. Total time decreases monotonically with base speed: from 1680.45 to 735.12~s (56.25\%) on \texttt{map\_test1} and from 14525.14 to 4120.11~s (71.63\%) on \texttt{parkingGraph}. As $\rho$ increases from 0.00556 to 0.03333~m/degree, coverage decreases from 96.44\% to 85.27\% and from 99.15\% to 90.32\%, reductions of 11.17 and 8.83 percentage points. Larger sweep spacing therefore weakens coverage continuity, but no tested condition makes the full-trajectory coverage unusably low.

Time per unit covered area further clarifies the trade-off. It decreases continuously with base speed on both maps: from 2.852 to 1.411~s/m$^2$ on \texttt{map\_test1} and from 3.412 to 1.062~s/m$^2$ on \texttt{parkingGraph}, improvements of 50.53\% and 68.88\%. A speed of 3.0~m/s gives both the shortest total time and lowest time per covered area, at the cost of coverage ratios of 85.27\% and 90.32\%. Conversely, 0.5~m/s gives the highest coverage on both maps. If both maps must retain at least 95\% full-trajectory coverage, 1.0~m/s is the highest tested admissible speed, while already reducing total time by 34.81\% and 42.54\% relative to 0.5~m/s. Base speed should therefore be selected jointly with arm angular velocity according to a minimum-coverage constraint and schedule objective, rather than treated as an incidental parameter after geometric planning. Smaller $\rho$ favors coverage quality; larger $\rho$ favors time efficiency when a coverage loss of roughly 9--11 percentage points is acceptable.

\begin{figure}[htbp]
  \centering
  \begin{minipage}[t]{0.49\linewidth}
    \centering
    \includegraphics[width=\linewidth]{figure_5a_vehicle_speed_map_simple_en.png}\\[-0.4em]
    \small (a) \texttt{map\_test1}
  \end{minipage}\hfill
  \begin{minipage}[t]{0.49\linewidth}
    \centering
    \includegraphics[width=\linewidth]{figure_5b_vehicle_speed_parkingGraph_en.png}\\[-0.4em]
    \small (b) \texttt{parkingGraph}
  \end{minipage}
  \caption{Effect of the base-to-arm velocity ratio on full-trajectory coverage and execution time. The horizontal axis is $\rho=v/\omega_a$; the secondary top axis gives the corresponding base speed. Blue denotes full-trajectory coverage and red total execution time. Both curves use the data in Table~\ref{tab:speed-study}.}
  \label{fig:speed-study}
\end{figure}

\subsection{Discussion}

The three experiments support one central conclusion: the contribution is not simply another decomposition, but the use of coordinated base motion and end-disc swinging throughout map decomposition, trajectory generation, and swept-area evaluation, so that coverage is determined by the disc's actual continuous motion. The Fields2Cover comparison shows that fixed-width strips are efficient in simple scenes, but coverage capability and connection feasibility can be constrained in dense obstacle fields; the time-varying end-disc sweep enlarges the effective workspace. The scan experiment shows that decomposition must reflect map structure and that bidirectional candidates reduce the representation bias of a single scan direction. The speed experiment shows that swing-arm CPP must jointly consider route geometry and execution speed, because matching base speed to arm angular velocity directly determines coverage quality.

These findings do not imply unconditional superiority over a fixed-width platform. In regular open areas, a strip-based method can obtain acceptable coverage with a shorter path and lower time. The value of the swing arm is greatest at boundaries a fixed base cannot closely follow, in obstacle-dense regions, and on wet concrete where fewer base maneuvers are desirable. Platform and parameter selection can therefore depend on site structure: large open regions favor efficient fixed-width sweeps, while complex corners and multi-obstacle regions favor arm-assisted coverage. The two could also operate in separate zones of the same floor.

The method has several limitations. First, rectangle candidates and visitation sequences are generated heuristically. They provide feasible and interpretable routes but no global optimality guarantee; many narrow corridors can still increase the rectangle count and repeated transit. Second, the input is a static two-dimensional map and does not explicitly account for workers, material stacks, localization drift, or changes in concrete condition during construction. Third, the speed analysis focuses on coverage density under base linear and arm angular velocities; it does not yet model full chassis dynamics, disc pressure, surface-quality feedback, or control error. Overall, the experiments identify explicit modeling of the end-tool sweep, area decomposition, and velocity coupling as the features that distinguish swing-arm troweling CPP from conventional fixed-width planning and define the next steps toward field deployment.
